\documentclass[10pt]{article}
\usepackage[margin=0.9in]{geometry}
\usepackage[T1]{fontenc}
\usepackage{lmodern}
\usepackage{amsmath,amssymb}
\usepackage{booktabs}
\usepackage{tabularx}
\usepackage{hyperref}
\usepackage{enumitem}
\usepackage{float}
\setlist{nosep}
\setlength{\parskip}{0.4em}
\setlength{\parindent}{0pt}
\newcommand{\den}[1]{[\![#1]\!]}
\newcommand{\plain}{\mathit{plain}}
\newcommand{\ceil}{P_{\mathrm{ceil}}}
\newcommand{\BC}{{\mdseries\textsc{BoundClaw}}}
\newcolumntype{L}[1]{>{\raggedright\arraybackslash}p{#1}}
\newcolumntype{Y}{>{\raggedright\arraybackslash}X}

\title{Policy Inference in \BC: Evaluation Methodology}
\date{Draft for advisor and collaborator review, \today}

\begin{document}
\maketitle

\section{What is evaluated}\label{sec:interface}
At each delegation, \BC{}'s policy engine asks an untrusted helper to select, from the delegator's
$(C_i, P_i)$ with $P = (N, \plain, D)$, the authority the subtask needs. \BC{}'s checker decides what is
allowed; the helper only decides how much of the least-privilege benefit a delegation gets. The checks in
our library (attenuation, ceiling, recovery admission) are advisory and exist so the benchmark can simulate
\BC{}'s checker offline. Table~\ref{tab:interface} fixes the interface every helper is evaluated under; it
follows the current draft and \BC{}'s adapter (\texttt{cmd/bc-mbla}). How each helper works, its exact prompts,
and the open input and output format questions are in the companion helper specification (\texttt{helper\_spec.pdf}).

\begin{table}[H]\centering\small
\begin{tabularx}{\linewidth}{L{0.15\linewidth}Y}
\toprule
Input & Root task (the user's request); subtask text written by the delegating agent; receiver identity and
kind (local or remote); current bounds $(C_i,P_i)$; upstream context; optional signed manifest, shown as a hint only. \\
Candidates & One per grant in $C_i$, plus a narrower grant for each resource the context names under a
granted prefix; one per destination in $N_i$ that has no parser (selected into $\plain'$); for a local receiver
only, one per volume mount $(v,\mu)$ in $D_i$. \emph{Not scored:} hosts with a parser (S3, httpapi) follow
$C$; peers are chosen by the agent's next-peer choice; the model endpoint is always kept. \\
Output & Candidate indices only, under a fixed JSON schema, in three tiers: required (the proposal),
plausible (the admission set $A$), unneeded. \\
Root-task wrapper & \BC{}'s adapter calls the helper twice in parallel, with the root task and with the
subtask, and keeps the lower score per candidate. Results are reported under this wrapper and on the subtask alone. \\
Failure & Timeout, error, unparsable output, an out-of-range index or an unknown key refuses the delegation.
There is no fallback proposal. \\
Receiver & $C$ is bounded by the parent only (no capability ceiling). A proposal whose $P'$ is not within the
receiver's $\ceil$ is rejected as a whole, not clipped; the delegator never sees $\ceil$. \\
Recovery & Started by the tunnel only. Each request names one atom or one host. It is granted only if it is
within the parent and in $A$ (a host without a parser must also be in $\plain_i$), and the receiver installs
a host only within $\ceil$. Volume denials are final. The cap $k$ counts requests, granted or not. \\
\bottomrule
\end{tabularx}
\caption{Helper interface. Every number in the paper is measured under these rules.}\label{tab:interface}
\end{table}

The evaluation does not need to show that inference is safe: G1 bounds every output, and a fully
compromised helper reproduces at most the parent's authority. It shows three things: (1) inference removes
authority that G1 alone would pass downstream without breaking tasks; (2) bounding recovery by $A$ keeps an
injected instruction from turning a denial into a grant; (3) the benefit survives an adversary who targets
the helper. Cost is the fourth axis, because the helper is on the path of every delegation.

\section{How the preliminary study shaped the design}\label{sec:prelim}
Before fixing this plan we ran an exploratory study: 23 helper configurations of four kinds on our
earlier policy model (findings report, F-01 to F-29). \textbf{Its numbers are not part of the final
evaluation.} We use them only to justify the choices below; every result in the paper comes from the
protocol in Sections~\ref{sec:rq}--\ref{sec:stats}.

\begin{table}[H]\centering\small
\begin{tabularx}{\linewidth}{L{0.33\linewidth}YL{0.07\linewidth}}
\toprule
Design choice & Preliminary evidence & Ref. \\
\midrule
Recovery bounded by $A$ is the main security question (RQ2) &
With parent-bounded recovery every planted attack stayed reachable; with $A$-bounded recovery 13\% did,
while 97\% of tasks finished & F-22, F-24 \\
The helper must emit three tiers &
Helpers with near-binary scores cannot separate the proposal from $A$; their tuned $A$ gave no protection & F-16 \\
Main candidates: mid-size generative helpers that see the whole list &
Qwen3-4B/8B/30B-A3B granted 6--9\% of unneeded authority; the largest model (gpt-oss-120b) 14\% & F-24 \\
Train fast helpers rather than drop them &
Qwen3-Reranker-0.6B decides in 28\,ms but grants 40\% of unneeded authority, against 51\% for a fair coin
& F-26, F-29 \\
Show the manifest as a hint; attack it with inflation &
The manifest cut unneeded authority from 23\% to 9\% for Qwen3-8B & F-24 \\
Compare with CAPMAS's method, not only its release &
As published (top-10 cut, no recovery) it granted 74\% of unneeded authority and finished 51\% of tasks & F-28 \\
Local, open-weight helpers only &
The helper reads possibly confidential context and its sandbox reaches only the model endpoint; open weights
also allow white-box attacks & F-09 \\
Whole volumes, no executable permissions &
The only reviewer disagreements were on path-level file and executable labels; the draft drops both & F-20 \\
\bottomrule
\end{tabularx}
\caption{Each design choice and the preliminary result that motivates it.}\label{tab:prelim}
\end{table}

\section{Research questions}\label{sec:rq}
\begin{description}
  \item[RQ1 Least privilege.] How much of the authority G1 alone would pass downstream does inference remove,
    per component ($C$, $N$, $\plain$, $D$), at what missing authority, and how often does over-selection get a
    proposal rejected at $\ceil$?
  \item[RQ2 Recovery.] Does bounding recovery by $A$ keep injected instructions from reaching authority outside
    the task, compared with a parent-bounded rule, and what does it cost legitimate tasks?
  \item[RQ3 Robustness.] How much of the RQ1 benefit survives planted injections, injection styles never seen in
    training, an adaptive attacker, a steered subtask (under \BC{}'s root-task wrapper and on the subtask
    alone), and an inflated manifest?
  \item[RQ4 Cost.] What latency does inference add per delegation (the wrapper's two calls run in parallel, so the slower one counts),
    and how often does the helper refuse?
  \item[RQ5 Training.] Does task-specific training let small, fast helpers match the best prompted helper on
    RQ1--RQ3, and does it hold on task families and attack styles never seen in training?
\end{description}

\section{Helpers}\label{sec:helpers}
\begin{table}[H]\centering\small
\begin{tabularx}{\linewidth}{L{0.25\linewidth}L{0.22\linewidth}Y}
\toprule
Helper & Variants & Role \\
\midrule
Qwen3-8B & prompted, fine-tuned & Best security--utility trade-off; the reference helper \\
Qwen3-4B, Qwen3-1.7B & prompted, fine-tuned & Smaller generative helpers \\
Qwen3-30B-A3B & prompted, QLoRA & Fastest precise generative helper; training must teach the plausible tier \\
Qwen3-Reranker-0.6B, -8B & prompted, fine-tuned & Fastest helper and strongest per-candidate classifier \\
CAPMAS~\cite{capmas} & as published; retrained on our data & Closest prior system \\
gpt-oss-120b & prompted & Upper reference point, not a deployment candidate \\
Set-only, manifest-only, lexical, coin flips (25/50/75\%) & -- & G1 alone, operator hint alone, keyword match, chance \\
\bottomrule
\end{tabularx}
\caption{Helpers. Others from the exploratory study were dominated by these or cannot be trained the same way.}
\end{table}

\section{Datasets}\label{sec:data}
\paragraph{Common format.} Every task is one delegation: root task, subtask, receiver (identity, local or
remote, $\ceil$), delegator holdings $(C_i, N_i, \plain_i, D_i)$, upstream context, optional manifest, and an
optional injected instruction with the authority it targets. Labels give, per component, the minimal
authority $X^\star$ (required) and the extra authority any alternative valid plan needs (plausible).
S3 atoms follow the prototype: (s3, verb, bucket, key prefix) with verb in \{GET, PUT, DELETE, LIST\}.
Hosts are parsed (S3, httpapi), plain, or peers. Volumes are $(v,\mu)$ with $\mu\in\{\mathrm{ro},\mathrm{rw}\}$.
There are no executable permissions.

\paragraph{Labeled evaluation set (primary; RQ1--RQ4).} 408 tasks from 17 templates over the paper's
workflows: 288 test and 120 development, with no template in both. Each template varies over 4 name
bindings, 3 levels of delegator over-provisioning (narrow, medium, broad: extra authority on the same
resource with another verb, on sibling resources, and wider prefixes), and with or without one injected
instruction aimed at authority the delegator holds but the task does not need. The injection sits in the
subtask or upstream context, never in the root task. Two reviewers label every candidate independently and
then adjudicate; we report Cohen's $\kappa$~\cite{cohen}. Test files are frozen by checksum before any helper
runs, and thresholds are tuned on development only.

\paragraph{Synthetic set (training and held-out generalization; RQ5).} No public dataset has our setting
(holdings as the choice set, ceilings, recovery, injections), so we generate one with methods from recent
least-privilege datasets~\cite{taskcond,toolpriv,authbench}.
\begin{itemize}
  \item \emph{Size and split.} 2,000 tasks in 40 families: 1,500 train and 300 validation from 32 families;
    200 test tasks from 8 families never trained on (seeded draw, stratified by workflow). Size follows the
    closest prior work~\cite{taskcond}. No name is shared with the labeled set, and tasks sharing any
    8-word sequence with it are dropped~\cite{contamination}.
  \item \emph{Generation.} A local open model writes the root task and subtask as the agents would, plus the
    API calls that complete the subtask and optionally an alternative valid plan~\cite{selfinstruct,apigen}.
    Receivers vary over local and remote, and tasks over one and two hops.
  \item \emph{Labels come from execution.} Each call runs against a mock service that checks its arguments
    and that every resource it touches is named in the context, and logs the atoms, hosts and volumes it uses.
    The required tier is the minimal candidate set covering the log, as IAM Access Analyzer derives policies
    from activity~\cite{accessanalyzer} and AuthBench derives labels from traces~\cite{authbench}; the
    plausible tier is what the alternative plan adds. Each call's permission comes from the provider's
    published reference (for S3, AWS's machine-readable service reference, checked automatically). The
    labeler reproduces the human labels on all 408 labeled tasks from their call plans.
  \item \emph{Checks.} Format, mock execution, and two judge models other than the writer, which must both
    agree the plan completes the task, uses no unneeded call and is unambiguous~\cite{toolpriv,apigen};
    near-duplicates are removed. We report the yield of each stage.
  \item \emph{Holdings and injections.} Built as in the labeled set; half the tasks carry an injection and
    its label does not change.
  \item \emph{Human acceptance test.} Two reviewers label 93 random tasks blind to the generated labels and
    adjudicate. The set is accepted if at most one generated label is wrong, which shows a label error rate
    below 5\% with 95\% confidence (one-sided binomial acceptance sampling~\cite{montgomery,clopper}).
    Otherwise we fix the generator, regenerate and draw a fresh sample.
\end{itemize}

\paragraph{External benchmarks.} FORTIS tasks 1 and 2~\cite{fortis} (official grading) and the mostargate
prompts~\cite{noyan}, to test whether conclusions transfer to another group's tasks and labels. Evaluation only.

\paragraph{End-to-end runs in \BC{} (RQ1, RQ2, RQ4).} The paper's object-storage workflows in \BC{}'s harness
(S3-compatible store, stub internet, scripted and live agents, deterministic oracle with canary strings), with
the subprocess proposer running each final helper through the adapter. We measure denials, recovery requests,
ceiling rejections, task completion, oracle-detected attack outcomes, and delegation latency through the tunnel.

\paragraph{Attack suite (RQ3).} On the labeled test set and the end-to-end runs: (a) the planted injections;
(b) injection styles from published corpora never used in training~\cite{injecagent,agentdojo}; (c) an adaptive
attacker that knows the design, candidate format and output schema, has the open weights, and searches for
text that places a target candidate in the proposal or in $A$ with an attacker-model loop under a fixed query
budget~\cite{carlini,attackersecond,pair}, reported as success versus budget; (d) a steered subtask, under the wrapper and on
the subtask alone; (e) a receiver operator who signs a manifest listing the attack's target as
usually needed.

\section{Training}\label{sec:training}
Training cannot add a guarantee; G1 supplies those. It changes precision and cost: the fast helpers failed on
task knowledge, not on the mechanism (Table~\ref{tab:prelim}). Training on injected examples with unchanged
labels also teaches a helper to ignore injected requests; held-out attack styles test whether that generalizes.
\begin{itemize}
  \item \emph{Targets.} Classifiers learn needed or not per candidate; generative helpers learn the three-tier
    index output; CAPMAS uses its own contrastive training code.
  \item \emph{Method.} Full fine-tuning of all weights for models up to 8B, so a trained model differs from its
    untrained version only by training; full fine-tuning learns new tasks better than low-rank
    adaptation~\cite{loralearnsless}. Qwen3-30B-A3B is too large for our hardware and uses QLoRA~\cite{qlora},
    compared only at matching precision.
  \item \emph{Settings.} Cosine schedule, 3\% warm-up, 2 epochs, loss on output tokens only; learning rate from
    $\{5\!\times\!10^{-6}, 10^{-5}, 2\!\times\!10^{-5}\}$ on synthetic validation; three seeds. Thresholds tuned on
    the labeled development split.
  \item \emph{Isolation.} No labeled, external, end-to-end or attack data is used for training or selection.
    The synthetic held-out families are used once, for the final report.
\end{itemize}

\section{Metrics and statistics}\label{sec:stats}
Per task and per component $X\in\{C,N,\plain,D\}$, against $X^\star$, comparing authority on witness atoms so a
prefix counts for everything it covers:
\begin{itemize}
  \item \textbf{Excess authority} $|\den{X'}\setminus\den{X^\star}| / |\den{X_i}\setminus\den{X^\star}|$ (set-only $=1$);
    \textbf{missing authority} $|\den{X^\star}\setminus\den{X'}| / |\den{X^\star}|$; exact match.
  \item \textbf{Ceiling rejections}: share of proposals rejected at $\ceil$.
  \item \textbf{Attack reachability after recovery}: share of attacks whose target is granted at delegation or
    reachable within $k=3$ requests under the rules of Table~\ref{tab:interface} (sensitivity $k\in\{1,5\}$),
    with $A$-bounded and parent-bounded admission.
  \item \textbf{Task completion} under the same rules (a missing volume, a refusal or a ceiling rejection is a
    failure); recovery requests per task; refusal rate.
  \item \textbf{Benefit retained under attack} $(E_{\mathrm{G1}}-E_{\mathrm{attack}})/(E_{\mathrm{G1}}-E_{\mathrm{clean}})$
    on the attacked authority: 1 means the attack gained nothing, 0 means the helper fell back to the parent.
  \item \textbf{Latency}: helper time (median, p95), each model alone on one GPU, and delegation setup time
    end to end.
\end{itemize}
\textbf{Protocol.} Helpers, metrics, comparisons and pass rules are fixed in our decision log before results
are seen; anything added later is labeled post hoc. Primary comparisons, each Holm-corrected within its
metric: best helper vs.\ set-only (RQ1); $A$-bounded vs.\ parent-bounded recovery (RQ2); each fine-tuned helper
vs.\ its untrained version and vs.\ prompted Qwen3-8B (RQ5); best helper vs.\ CAPMAS and vs.\ each coin flip.
Confidence intervals (95\%) come from bootstrap resampling of templates or task families, since tasks within one
are correlated; paired binary outcomes use McNemar's exact test; trained models report the mean and spread over
seeds.

\section{Open decisions and the \BC{} asks}\label{sec:decisions}
\begin{enumerate}
  \item \textbf{Services beyond S3.} Only S3 and httpapi have parsers. We propose keeping the issue-tracker,
    document and chat workflows as httpapi atoms if httpapi's atom form fits them; otherwise their hosts are
    plain only and the capability results come from S3 alone.
  \item Local or remote receiver for each handoff, and $\ceil$ for each receiver role.
  \item The cap $k$ (we propose 3) and the helper timeout (prototype: 60\,s).
  \item Which hosts are plain in each workflow.
  \item Whether end-to-end runs may log helper candidates, tiers and $A$.
\end{enumerate}
Input and output format questions on the model side (prompt structure, candidate order, training targets) are listed
with recommendations in the helper specification and are decided by our side before the final runs.

\begin{table}[H]\centering\small
\begin{tabularx}{\linewidth}{L{0.06\linewidth}Y}
\toprule
Ask & Where this plan meets it \\
\midrule
B1 & $P=(N,\plain,D)$, whole volumes, no executables (\S\ref{sec:data}, common format) \\
B2 & Volumes granted at delegation; volume denials final (Table~\ref{tab:interface}) \\
B3 & Note from \BC{}: its adapter applies the root-task rule. Every task carries a root task and results are reported under it (RQ3d) \\
B4 & Hosts are parsed, plain or peers; only unparsed hosts are scored (Table~\ref{tab:interface}) \\
B5 & $C$ bounded by the parent; ceiling miss rejects the delegation and is a reported metric \\
B6 & $k$ counts requests; fixed in \texttt{Recovery.Decide} \\
L1 & Out-of-range indices, unknown keys and trailing text refuse; fixed in the chat scorer \\
D1--D3 & Library README and spec now say \BC{}'s checker decides, describe recovery as above, and reject at the ceiling \\
D4 & Paper-alignment notes remapped to the current draft (with the benchmark rebuild) \\
\bottomrule
\end{tabularx}
\caption{How the plan answers \BC{}'s requests to the policy-inference side (Oct 26).}
\end{table}

\begin{thebibliography}{99}\small
\bibitem{capmas} R.~M.~Veski, R.~Guerraoui, D.~Froelicher. CAPMAS: Capability-based delegation of privileges in multi-agent systems. arXiv:2609.06500, 2026.
\bibitem{cohen} J.~Cohen. A coefficient of agreement for nominal scales. Educational and Psychological Measurement, 1960.
\bibitem{taskcond} Task-conditioned least-privilege learning for executable terminal and MCP agents. arXiv:2608.18351, 2026.
\bibitem{toolpriv} K.~Yang et al. When lower privileges suffice: Investigating over-privileged tool selection in LLM agents. arXiv:2606.20023, 2026.
\bibitem{authbench} Do coding agents understand least-privilege authorization? (AuthBench). arXiv:2605.14859, 2026.
\bibitem{contamination} O.~Sainz et al. NLP evaluation in trouble: On the need to measure LLM data contamination for each benchmark. Findings of EMNLP, 2023.
\bibitem{selfinstruct} Y.~Wang et al. Self-Instruct: Aligning language models with self-generated instructions. ACL, 2023.
\bibitem{apigen} Z.~Liu et al. APIGen: Automated pipeline for generating verifiable and diverse function-calling datasets. arXiv:2406.18518, 2024.
\bibitem{accessanalyzer} Amazon Web Services. IAM Access Analyzer policy generation. AWS documentation.
\bibitem{montgomery} D.~C.~Montgomery. \emph{Introduction to Statistical Quality Control}, chapter on lot-by-lot acceptance sampling. Wiley.
\bibitem{clopper} C.~J.~Clopper, E.~S.~Pearson. The use of confidence or fiducial limits illustrated in the case of the binomial. Biometrika, 1934.
\bibitem{fortis} FORTIS benchmark for least-privilege skill and tool selection. \url{https://github.com/lili0415/FORTIS-Benchmark}.
\bibitem{noyan} H.~B.~Noyan. Dynamic capability scoping for enterprise AI agents: A synthetic dataset and three-source permission architecture. ICML 2026; arXiv:2607.22445.
\bibitem{injecagent} Q.~Zhan, Z.~Liang, Z.~Ying, D.~Kang. InjecAgent: Benchmarking indirect prompt injections in tool-integrated large language model agents. 2024.
\bibitem{agentdojo} E.~Debenedetti et al. AgentDojo: A dynamic environment to evaluate prompt injection attacks and defenses for LLM agents. NeurIPS Datasets and Benchmarks, 2024.
\bibitem{carlini} N.~Carlini et al. On evaluating adversarial robustness. arXiv:1902.06705, 2019.
\bibitem{attackersecond} M.~Nasr et al. The attacker moves second: Stronger adaptive attacks bypass defenses against LLM jailbreaks and prompt injections. arXiv:2510.09023, 2025.
\bibitem{pair} P.~Chao et al. Jailbreaking black box large language models in twenty queries. arXiv:2310.08419, 2023.
\bibitem{loralearnsless} D.~Biderman et al. LoRA learns less and forgets less. TMLR, 2024.
\bibitem{qlora} T.~Dettmers et al. QLoRA: Efficient finetuning of quantized LLMs. NeurIPS, 2023.
\end{thebibliography}
\end{document}
